% =====================================================================
% EE 311 Project Plan (IEEE journal format)
% Compile with pdflatex, or upload this single file to Overleaf.
% =====================================================================
\documentclass[journal]{IEEEtran}

\usepackage[T1]{fontenc}
\usepackage[utf8]{inputenc}
\usepackage{cite}
\usepackage{booktabs}
\usepackage{tabularx}
\usepackage{url}

% ---- Author (single-person project) ----
\newcommand{\authorname}{Oğuz Kağan Göçmez}

\begin{document}

\title{Machine Learning Prediction of Heating and Cooling Loads of Residential Buildings\\[0.3em]
{\LARGE Project Plan}}

\author{\authorname%
\thanks{The author is with Özyeğin University, İstanbul, Türkiye.
This plan is part of the EE~311 Introduction to Artificial Intelligence,
Machine Learning and Data Science course project, Fall 2026.}}

\maketitle

% ---------------------------------------------------------------------
\begin{abstract}
Heating and cooling loads determine the size of the heating, ventilation and air conditioning (HVAC)
equipment that a building needs. They are usually computed with detailed simulation software, which is
slow and requires expert users. Machine learning (ML) models trained on simulated buildings can give
the same estimates in a fraction of a second. In this project we use the UCI Energy Efficiency dataset,
which contains 768 residential buildings simulated in Ecotect and described by eight design parameters
(relative compactness, surface area, wall area, roof area, overall height, orientation, glazing area and
glazing area distribution). We predict two outputs: the heating load (HL) and the cooling load (CL).
Following the reference study of Tsanas and Xifara, we will explore the data statistically, compare six
regression models with repeated 10-fold cross-validation and three error measures, tune the most
promising models, and analyze which inputs matter most.

As our own contribution, we will test how well the models generalize to building shapes that were not
seen during training. In a random split, buildings that differ only in orientation, which mostly have
very similar loads, end up in both the training and the test sets, which may make the reported errors
optimistic. Our first experiments already reproduce the statistical analysis of the reference study, and
an untuned random forest reaches a mean absolute error of 0.31 for HL and 0.98 for CL, compared with
0.51 and 1.42 in the reference study.
\end{abstract}

\begin{IEEEkeywords}
Building energy performance, heating load, cooling load, machine learning, regression, cross-validation.
\end{IEEEkeywords}

% ---------------------------------------------------------------------
\section{Introduction}
\IEEEPARstart{H}{VAC} systems account for most of the energy used in buildings, so estimating heating
and cooling loads early in the design phase helps to design more energy-efficient buildings~\cite{tsanas2012}.
Simulation tools give reliable estimates but are time-consuming, which has motivated the use of ML
models such as decision trees~\cite{yu2010}, support vector machines~\cite{dong2005} and polynomial
regression~\cite{catalina2008}.

Tsanas and Xifara~\cite{tsanas2012} published the dataset used in this project~\cite{uci2012}. They
compared a robust linear model, iteratively reweighted least squares (IRLS), with a random forest (RF),
and reported mean absolute errors (MAE) of 0.51 for HL and 1.42 for CL with the RF. Later studies
applied other artificial intelligence techniques~\cite{chou2014} and genetic programming~\cite{castelli2015}
to the same data. The dataset is also a common benchmark: \v{Z}egklitz and Po\v{s}\'{\i}k~\cite{zegklitz2017}
used it to compare symbolic regression methods with a tuned RF and support vector regression (SVR), and
Yazici \emph{et al.}~\cite{yazici2018} used it to test ML models on a low-power Internet of Things device.
We follow the Machine Learning Project Checklist of the course textbook~\cite{geron2022}.

% ---------------------------------------------------------------------
\section{Project Schedule}
Table~\ref{tab:schedule} lists the subtasks with their deadlines. This is a single-person project, so
the author is responsible for all subtasks. The approach for each subtask is described below.

\begin{table}[!t]
\caption{Project Schedule}
\label{tab:schedule}
\centering
\begin{tabularx}{\columnwidth}{@{}Xcc@{}}
\toprule
Subtask & Week & Status \\
\midrule
Literature search               & 1--2   & In progress \\
Frame the problem               & 1      & Done \\
Get the data                    & 1      & Done \\
Explore the data                & 2      & Done \\
\textbf{Project plan report}    & 2      & This report \\
Prepare the data                & 3      & Planned \\
Shortlist models (random CV)    & 3--4   & First pass \\
Grouped CV by building shape    & 5      & Planned \\
\textbf{Midterm report}         & 6      & Planned \\
Fine-tune the best models       & 7--8   & Planned \\
Feature importance and errors   & 9      & Planned \\
Final tables and figures        & 10     & Planned \\
\textbf{Final report}           & 11--12 & Planned \\
\textbf{Recorded presentation}  & 12     & Planned \\
\bottomrule
\end{tabularx}

\vspace{2pt}
\footnotesize Week 1 is the week the topic was selected; week 12 is the last week of the semester.
All subtasks are carried out by the author.
\end{table}

\subsection{Literature Search}
We read the reference study~\cite{tsanas2012} and its references, and we search for later work on the
same dataset (for example, papers citing~\cite{tsanas2012} on Google Scholar) to collect results to
compare against. This continues until the final report.

\subsection{Frame the Problem}
The task is supervised regression with two continuous outputs, HL and CL. To make our results directly
comparable with~\cite{tsanas2012}, we use the same three errors: mean absolute error (MAE), mean squared
error (MSE) and mean relative error (MRE). The IRLS and RF results of~\cite{tsanas2012} are our baseline.

\subsection{Get the Data}
We load the dataset~\cite{uci2012} from a CSV export of the original Excel file, which uses semicolons
as separators and commas as decimal points. We checked it against the Excel file: 768 buildings,
10 columns, no missing values and no duplicate rows.

\subsection{Explore the Data}
We plot histograms and input--output scatter plots, and we measure the association of each input with
each output using the Spearman rank correlation (with $p$-values) and mutual information, as
in~\cite{tsanas2012}. This step is done, and our Spearman correlations match those of~\cite{tsanas2012}.
Relative compactness and surface area are perfectly rank-correlated ($\rho=-1$), so they carry the same
information. Orientation is not correlated with either output, although it does change the cooling load
of some two-storey shapes, in a way that does not go in one direction.

\subsection{Prepare the Data}
No cleaning is needed. Inputs are standardized inside a pipeline, so the scaler is fitted on the
training folds only. We will also test one-hot encoding of orientation and glazing area distribution,
and removing the redundant surface area input.

\subsection{Shortlist Promising Models}
We compare six models implemented in scikit-learn~\cite{pedregosa2011}: linear regression, robust
linear regression (Huber loss, our counterpart to IRLS), $k$-nearest neighbors, support vector regression,
random forest~\cite{breiman2001} and gradient boosting~\cite{friedman2001}. They are evaluated with
10-fold cross-validation (CV) repeated on random splits, as in~\cite{tsanas2012}. As our contribution,
we repeat the evaluation with grouped CV~\cite{roberts2017}, where all buildings of the same shape (one of
12) stay in the same fold, to measure performance on unseen shapes. Table~\ref{tab:preliminary} shows the
first, untuned results.

\begin{table}[!t]
\caption{First Results: MAE (Mean $\pm$ Std) With 10-Fold CV Repeated 10 Times, Untuned Models}
\label{tab:preliminary}
\centering
\begin{tabular}{@{}lcc@{}}
\toprule
Model & Heating load & Cooling load \\
\midrule
Linear regression           & $2.09 \pm 0.24$ & $2.27 \pm 0.27$ \\
Robust linear (Huber)       & $2.03 \pm 0.25$ & $2.20 \pm 0.30$ \\
$k$-nearest neighbors       & $1.42 \pm 0.20$ & $1.66 \pm 0.24$ \\
Support vector regression   & $0.55 \pm 0.08$ & $1.13 \pm 0.19$ \\
Gradient boosting           & $0.34 \pm 0.04$ & $1.02 \pm 0.13$ \\
Random forest               & $\mathbf{0.31 \pm 0.04}$ & $\mathbf{0.98 \pm 0.14}$ \\
\midrule
IRLS in~\cite{tsanas2012}   & $2.14 \pm 0.24$ & $2.21 \pm 0.28$ \\
RF in~\cite{tsanas2012}     & $0.51 \pm 0.11$ & $1.42 \pm 0.25$ \\
\bottomrule
\end{tabular}
\end{table}

\subsection{Fine-Tune the System}
For the two or three best models we will search the hyperparameters (grid or randomized search inside
CV). The final results will use 100 CV repetitions, as in~\cite{tsanas2012}. We will also compute
permutation importance and add a random probe feature, as in~\cite{tsanas2012}, to check which inputs
matter most, and we will try a single model that predicts both outputs together.

\subsection{Present the Results}
We will summarize the results in tables and figures, compare them with~\cite{tsanas2012},
\cite{chou2014}, \cite{castelli2015} and~\cite{zegklitz2017}, and prepare the midterm report, the final report and a recorded
presentation of 10--15 minutes.

% ---------------------------------------------------------------------
\section*{Acknowledgment}
AI tools (Claude, Anthropic) were used to help draft this plan and the initial analysis code.
All suggestions were reviewed carefully by the author, and their correctness was validated.

% ---------------------------------------------------------------------
\begin{thebibliography}{12}

\bibitem{tsanas2012}
A.~Tsanas and A.~Xifara, ``Accurate quantitative estimation of energy performance of residential
buildings using statistical machine learning tools,'' \emph{Energy Build.}, vol.~49, pp.~560--567,
2012, doi: 10.1016/j.enbuild.2012.03.003.

\bibitem{yu2010}
Z.~Yu, F.~Haghighat, B.~C.~M. Fung, and H.~Yoshino, ``A decision tree method for building energy
demand modeling,'' \emph{Energy Build.}, vol.~42, no.~10, pp.~1637--1646, 2010,
doi: 10.1016/j.enbuild.2010.04.006.

\bibitem{dong2005}
B.~Dong, C.~Cao, and S.~E. Lee, ``Applying support vector machines to predict building energy
consumption in tropical region,'' \emph{Energy Build.}, vol.~37, no.~5, pp.~545--553, 2005,
doi: 10.1016/j.enbuild.2004.09.009.

\bibitem{catalina2008}
T.~Catalina, J.~Virgone, and E.~Blanco, ``Development and validation of regression models to predict
monthly heating demand for residential buildings,'' \emph{Energy Build.}, vol.~40, no.~10,
pp.~1825--1832, 2008, doi: 10.1016/j.enbuild.2008.04.001.

\bibitem{uci2012}
A.~Tsanas and A.~Xifara, ``Energy Efficiency,'' UCI Machine Learning Repository, 2012.
[Online]. Available: \url{https://doi.org/10.24432/C51307}

\bibitem{chou2014}
J.-S. Chou and D.-K. Bui, ``Modeling heating and cooling loads by artificial intelligence for
energy-efficient building design,'' \emph{Energy Build.}, vol.~82, pp.~437--446, 2014,
doi: 10.1016/j.enbuild.2014.07.036.

\bibitem{castelli2015}
M.~Castelli, L.~Trujillo, L.~Vanneschi, and A.~Popovi\v{c}, ``Prediction of energy performance of
residential buildings: A genetic programming approach,'' \emph{Energy Build.}, vol.~102, pp.~67--74,
2015, doi: 10.1016/j.enbuild.2015.05.013.

\bibitem{zegklitz2017}
J.~\v{Z}egklitz and P.~Po\v{s}\'{\i}k, ``Symbolic regression algorithms with built-in linear
regression,'' 2017, arXiv:1701.03641. [Online]. Available: \url{https://arxiv.org/abs/1701.03641}

\bibitem{yazici2018}
M.~T. Yazici, S.~Basurra, and M.~M. Gaber, ``Edge machine learning: Enabling smart Internet of Things
applications,'' \emph{Big Data Cogn. Comput.}, vol.~2, no.~3, Art. no.~26, 2018,
doi: 10.3390/bdcc2030026.

\bibitem{geron2022}
A.~G\'{e}ron, \emph{Hands-On Machine Learning with Scikit-Learn, Keras, and TensorFlow}, 3rd~ed.
Sebastopol, CA, USA: O'Reilly Media, 2022.

\bibitem{pedregosa2011}
F.~Pedregosa \emph{et al.}, ``Scikit-learn: Machine learning in Python,'' \emph{J. Mach. Learn. Res.},
vol.~12, pp.~2825--2830, 2011. [Online]. Available: \url{https://jmlr.org/papers/v12/pedregosa11a.html}

\bibitem{breiman2001}
L.~Breiman, ``Random forests,'' \emph{Mach. Learn.}, vol.~45, no.~1, pp.~5--32, 2001,
doi: 10.1023/A:1010933404324.

\bibitem{friedman2001}
J.~H. Friedman, ``Greedy function approximation: A gradient boosting machine,''
\emph{Ann. Statist.}, vol.~29, no.~5, pp.~1189--1232, 2001, doi: 10.1214/aos/1013203451.

\bibitem{roberts2017}
D.~R. Roberts \emph{et al.}, ``Cross-validation strategies for data with temporal, spatial,
hierarchical, or phylogenetic structure,'' \emph{Ecography}, vol.~40, no.~8, pp.~913--929, 2017,
doi: 10.1111/ecog.02881.

\end{thebibliography}

\end{document}
